\chapter{Basic concepts of enriched category theory}

There are many reasons to consider enrichments, but one notable one for this text is the desire for a ``homotopy product'', which would be a product diagram that only commutes \emph{up to homotopy}. In order to have some notion of homotopy for which this makes sense, we can assume hom-sets are equipped with topologies, and then homotopies between arrows would become paths in the underlying space. A homotopy product can then be equipped with a natural weak homotopy equivalence.

Since we have not defined a class of weak equivalences, the homotopy category $h\mathcal{M} $ is different from the ones we worked with earlier. Denoting by $\underline{\mathcal{M} }(a,b)$ the space assigned to $\mathcal{M} (a,b)$, we define $h\mathcal{M} (a,b)\coloneqq \pi_0 \underline{\mathcal{M} }(a,b)$. Here note that $\pi_0$ is the functor $\Top \to \Set $ assigning to each space it's path components; hence taking a space $X$ and quotienting in $\Set $ via the relation $x\sim y$ iff there is a path from $x$ to $y$.

\section{A first example}

Let $\Rmod_R$ be the category of left $R$-modules. We well know that the set $\Rmod_R(A,B)$ admits an abelian group structure via pointwise addition, which we will denote by $\underline{\Rmod_R}(A,B)$. Moreover, composition
\[
	\circ :\Rmod_R(B,C)\times \Rmod_R(A,B) \to \Rmod_R(A,C)
\]
distributes over addition, and hence is $\mathbb{Z} $-bilinear. Thus it factors through $\otimes _\mathbb{Z} $, inducing an abelian group homomorphism
\[
	\underline{\Rmod_R}(B,C)\otimes_\mathbb{Z}  \underline{\Rmod_R}(A,B) \to \underline{\Rmod_R}(A,C)
\]
in the category $\Ab$.

Note also that since any abelian group $a$ has $\Ab(\mathbb{Z} ,a)\cong a$, we can represent the identity morphism of $A$ as an abelian group homomorphism
\[
	1_A : \mathbb{Z}  \to \underline{\Rmod}_R(A,A)
.\]
Moreover, $\mathbb{Z} $ is a unit for the monoidal product $\otimes _\mathbb{Z} : \Ab\times \Ab \to \Ab$ via left and right natural isomorphisms. This allopws us to encode the usual composition axioms via diagrams in $\Ab$. These facts show that $\Rmod_R$ admits the structure of a category enriched over $(\Ab,\otimes _\mathbb{Z} ,\mathbb{Z} )$.

\section{The base for enrichment}

A symmetric monoidal category $(\mathcal{V} ,\times ,*)$ is a category with a bifunctor $-\times -: \mathcal{V} \times \mathcal{V}  \to \mathcal{V} $ called the monoidal product, and an element $*\in \mathcal{V} $ called the unit, together with specified natural isomorphisms
\[
	v\times w\cong w\times v,\qquad u\times (v\times w)\cong (u\times v)\times w,\qquad *\times v\cong v\cong v\times *
\]
satisfying coherence criteria. Any category with finite products is symmetric monoidal with terminal object the unit; we call a category with such a monoidal structure \emph{cartesian monoidal}. In most relevant examples, $\mathcal{V} $ is complete and cocomplete, so we will generally assume this.

\section{Enriched categories}

\begin{definition}[Enriched Category]\label{dfn:3.1}
	Let $\mathcal{V} $ be a symmetric\footnote{We are using symmetric monoidal categories here, but the symmetric hypothesis can be dropped for general enrichments. We require it for some results, however. I will assume we only enrich over symmetric monoidal categories unless the text specifies otherwise.} monoidal category. A $\mathcal{V} $-category $\underline{\mathcal{D} }$ is
	\begin{itemize}
		\item A collection of objects $x,y,z\in \underline{\mathcal{D}} $;
		\item For each pair $x,y\in \underline{\mathcal{D} }$, a hom-\emph{object} $\underline{\mathcal{D} }(x,y)\in \mathcal{V} $;
		\item For each $x\in\underline{\mathcal{D} }$, a morphism $\mathrm{id}_x : * \to \underline{\mathcal{D} }(x,x)$ in $\mathcal{V} $;
		\item For each triple $x,y,z\in\underline{\mathcal{D} }$, a morphism $\circ : \underline{\mathcal{D} }(y,z)\otimes \underline{\mathcal{D} }(xy) \to \underline{\mathcal{D} }(x,z)$ in $\mathcal{V} $;
	\end{itemize}
	such that for all $w,x,y,z\in\underline{\mathcal{D} }$, the diagrams
	\[\begin{tikzcd}[row sep = 3.5em, column sep = 3.5em]
		\underline{\mathcal{D} }(z,y)\otimes \underline{\mathcal{D} }(x,y)\otimes \underline{\mathcal{D} }(w,x)\ar[" 1\otimes \circ  ", r] \ar[" \circ \otimes 1 "', d]  & \underline{\mathcal{D} }(z,y)\otimes \underline{\mathcal{D} }(x,z)\ar[" \circ  ", d]  \\
		\underline{\mathcal{D} }(y,w)\otimes \underline{\mathcal{D} }(w,x)\ar[" \circ  "', r]  & \underline{\mathcal{D} }(w,z)
	\end{tikzcd}\]
	\[\begin{tikzcd}[row sep = 3.5em, column sep = 3.5em]
		\underline{\mathcal{D} }(x,y)\otimes *\ar[" 1\otimes \mathrm{id}_x ", r] \ar[" \cong "', dr]  & \underline{\mathcal{D} }(x,y)\otimes \underline{\mathcal{D} }(x,x)\ar[" \circ  ", d]  & \underline{\mathcal{D} }(y,y)\otimes \underline{\mathcal{D} }(x,y)\ar[" \circ  ", d]  & *\otimes \underline{\mathcal{D} }(x,y)\ar[" \mathrm{id} _y \otimes 1 "', l] \ar[" \cong ", dl]  \\
													      & \underline{\mathcal{D} }(x,y) & \underline{\mathcal{D} }(x,y) & 
	\end{tikzcd}\]
	commute in $\mathcal{V}$.
\end{definition}

We omit a relevant associativity morphism in the first diagram because the text did it too and I only realized after copying it all down and I'm NOT editing it.

We already know of a plethora of examples. For example, a ring is an $\Ab$-category with a singleton for objects. A more interesting example is that topological spaces are \emph{naturally} enriched over groupoids. Define $\underline{\Top }(X,Y)$ to be the groupoid with objects as continuous maps $X\to Y$ and morphisms as homotopy classes of homotopies between these maps.

If $\mathcal{V} $ has copowers of the unit that are preserved by tensor in the following sense, then any unenriched $\mathcal{C} $ has a free $\mathcal{V} $-category with objects the same as $\mathcal{C} $. Hom-objects from $a$ to $b$ are given by the copower $\coprod_{\mathcal{C} (a,b)} *$. Identity morphisms $\mathrm{id} _a$ are given by the inclusion of $*$ into the copower $\coprod_{\mathcal{C} (a,a)} *$ at the identity $1_a$. Since tensor preserves the relevant copower, we obtain
\[
	\left( \coprod_{\mathcal{C} (b,c)} *  \right) \otimes \left( \coprod_{\mathcal{C} (a,b)} * \right) \cong \coprod_{\mathcal{C} (b,c)} \left( *\otimes \left( \coprod_{\mathcal{C} (a,b)} * \right)  \right) \cong \coprod_{\mathcal{C} (b,c)} \left( \coprod_{\mathcal{C} (a,b)} \left( *\otimes * \right)  \right) \cong \coprod_{\mathcal{C} (b,c)} \coprod_{\mathcal{C} (a,b)} *\cong \coprod_{\mathcal{C} (b,c)\times \mathcal{C} (a,b)} *
.\]
Since there is a composition function $\mathcal{C} (b,c)\times \mathcal{C} (a,b)$, composition in the free $\mathcal{V} $-category simply reindexes along this arrow.

\begin{definition}[Closed Monoidal Category]\label{dfn:3.2}
	Let $\mathcal{V} $ be a symmetric monoidal category. If each functor $-\otimes v : \mathcal{V}  \to \mathcal{V} $ has a right adjoint, denoted $\underline{\mathcal{V} }(v,-)$, the right adjoints assemble uniquely into a bifunctor
	\[
		\underline{\mathcal{V} }(-,-): \mathcal{V} ^{\mathrm{op}} \times  \mathcal{V}  \to \mathcal{V} 
	\]
	such that for all $u,v,w\in \mathcal{V} $,
	\[
		\mathcal{V} (u\times v,w)\cong \mathcal{V} (u,\underline{\mathcal{V} }(v,w))
	\]
	natural in $u,v,w$.

	We call $\underline{\mathcal{V} }(-,-)$ the \emph{internal hom}, and $\mathcal{V} $ is enriched over itself with composition
	\[
		\underline{\mathcal{V} }(v,w)\otimes \underline{\mathcal{V} }(u,v)\to \underline{\mathcal{V} }(u,w)
	\]
	given as follows. Note the adjunction $-\otimes u\dashv \underline{\mathcal{V} }(u,-)$ gives the counit
	\[
		\varepsilon^u : (-\otimes u)\circ \underline{\mathcal{V} }(u,-)=\underline{\mathcal{V} }(u,-)\otimes u \Rightarrow 1
	.\]
	Thus the adjunctions give a composition
	\[\begin{tikzcd}[row sep = 2.5em, column sep = 2.5em]
		\underline{\mathcal{V} }(v,w)\otimes \underline{\mathcal{V} }(u,v)\otimes u\ar[" 1\otimes \varepsilon _v^u ", r]  & \underline{\mathcal{V} }(v,w)\otimes v\ar[" \varepsilon _w^v ", r]  & w.
	\end{tikzcd}\]
	Hence we have
	\[
		\varepsilon _w^v\circ (\varepsilon _v^u \otimes 1)\in \mathcal{V} \left( \underline{\mathcal{V} }(v,w)\otimes \underline{\mathcal{V} }(u,v)\otimes u,w \right) \cong \mathcal{V} \left( \underline{\mathcal{V} }(v,w)\otimes \underline{\mathcal{V} }(u,v),\underline{\mathcal{V} }(u,w) \right) 
	.\]
	The isomorphism follows by applying symmetry to the adjunction. We thus define composition as the adjunct of the relevant composition. The identity morphisms are defined as follows. The natural isomorphism $\lambda_v : *\otimes v \cong v$ is adjunct under the same adjunction to the identity $\mathrm{id}_v$:
	\[
		\lambda\in \mathcal{V} (*\otimes v,v)\cong \mathcal{V} (*,\underline{\mathcal{V} }(v,v))
	.\]
	Such a triple $(\mathcal{V} ,\otimes,*)$ is a \emph{closed symmetric monoidal category}.
\end{definition}

When $\mathcal{V} $ is cartesian monoidal, we call it \emph{cartesian closed}. An example is $\sSet$. The product is the usual product of functors, and we desire an adjunction
\[
	\sSet(X\times Y,Z)\cong\sSet(X,Z^Y)
.\]
Here, we use exponential notation $\underline{\sSet}(X,Y)\eqqcolon Y^X$. If $X$ is representable $\boldsymbol{\Delta} (-,[n])$, which we write as $\Delta^n$, then Yoneda gives
\[
	(Z^Y)_n\cong\sSet(\Delta^n, Z^Y)\cong\sSet(\Delta^n\times Y,Z)
.\]
Hence we obtain a valid definition of $Z^Y$ by setting the $n$-simplices as the set of maps $\Delta^n\times Y\to Z$. Hence $-\times Y\dashv (-)^Y$ for each simplicial set $Y$, giving that $\sSet$ is a closed cartesian monoidal category.

The category $\Cat $ is cartesian closed, since functors and natural transformations between 2 fixed categories form another category. The precise definition of a 2-category is a category enriched over $\Cat $. Objects are 0-cells, the 1-cells are the objects of the hom-categories, and the 2-cells are the morphisms in the hom-categories. This interpretation gives us a notion of arrows in the enriched category, compatable with the notion of composition given in \cref{dfn:3.2}. Equivalently, we can take 1-cells to be the arrows of the \emph{underlying category}, defined later.

The category $\Top $ is \emph{not} cartesian closed, but several large convenient subcategories are, such as $\mathcal{C} \mathcal{G} \mathcal{H} \mathrm{aus} $. For now, $\Top $ will denote one of these convenient categories, and we will make precise which categories we concern ourselves with later in the text.

We provide the following construction. Let $(\mathcal{V} ,\times ,*)$ be a cartesian closed symmetric monoidal category, which we also assume to be complete and cocomplete. There is a general procedure for producing a closed symmetric monoidal category $(\mathcal{V} _*,\wedge,S^0)$ of \emph{pointed objects} or \emph{based objects}, equivalently the slice category $(*\downarrow \mathcal{V} )\eqqcolon \mathcal{V} _*$. We define $S^0\coloneqq *\amalg *$ (the coproduct in $\mathcal{V}$). Note that for each $v,w$, we can define a map $v\to w$ by factoring through the terminal object $v\to *\to w$ via the basepoint on $w$. This gives maps $v,w\to v\times w$, and universality of coproduct gives a canonical map $f:v\amalg w \to v\times w$.We then define the \emph{smash product} $\wedge$ to be the pushout in $\mathcal{V} $:
\[\begin{tikzcd}[row sep = 2.5em, column sep = 2.5em]
v\amalg w\ar["f", r] \ar[d]  & v\times w\ar[d] \ar[bend left, ddr]  &  \\
*\ar[r] \ar[bend right, drr]  & v\wedge w \ar[dashed, dr] &  \\
 &  & z.
\end{tikzcd}\]
We then define the hom-objects $\underline{\mathcal{V} }_*$ in $\mathcal{V} _*$ to be pullbacks of internal homs in $\mathcal{V} $:
\[\begin{tikzcd}[row sep = 2.5em, column sep = 2.5em]
z\ar[bend left, drr] \ar[bend right, ddr] \ar[dashed, dr]  &  &  \\
							   & \underline{\mathcal{V} }_*(v,w)\ar[r] \ar[d]  & \underline{\mathcal{V} }(v,w)\ar[d]  \\
							   & \underline{\mathcal{V} }(*,*)\ar[r]  & \underline{\mathcal{V} }(*,w).
\end{tikzcd}\]
The bottom map is postcomposition by the base map $*\to w$, and the downwards map is precomposition by the basepoint $*\to v$. This does not have a precise definition yet, but intuitively it makes sense, since hom-objects in cartesian monoidal categories are intuitively similar to hom-sets. There is then a forgetful functor $U : \mathcal{V} _* \to \mathcal{V} $ forgetting the basepoint, which is left adjoint to the functor $(-)_+ : \mathcal{V}  \to \mathcal{V}_+ $ given by computing the coproduct with the terminal object.

\begin{lemma}\label{lma:3.1}
	When $\mathcal{V} $ is a cartesian closed symmetric monoidal category, the functor $(-)_+ : \mathcal{V}  \to \mathcal{V} _*$ is \emph{strong monoidal}, meaning
	\[
		S^0\cong(*)_+,\qquad v_+\wedge w_+\cong(v\times w)_+
	\]
	with these isomorphisms natural, and appropriately commuting with associators and unitors.
\end{lemma}
\begin{proof}
	The unit isomorphism is by definition of $S^0$. By definition, the object $v_+\wedge w_+$ is given by the cofiber (the opposite arrow in the pushforward diagram) of
	\[
	v_+\amalg w_+ \to v_+\times w_+=(v\amalg *)\times (w\amalg *)
	.\]
	Since $\mathcal{V} $ is cartesian closed and symmetric, $v\times -$ and $-\times v$ are left-adjoint and thus preserve colimits, and thus preserve coproducts. Thus, $\times $ distributes over $\amalg,$ and we obtain
	\[
		(v\amalg *)\times (w\amalg *)\cong (v\times (w\amalg *))\amalg(*\times(w\amalg*))\cong((v\times w)\amalg(v\times *))\amalg((*\times w)\amalg(*\times *))
	.\]
	Dropping the unnecessary parenthases due to associativity of $\amalg$, we also have
	\[
		(v\times w)\amalg(v\times *)\amalg(*\times w)\amalg(*\times *)\cong(v\times w)\amalg v\amalg w\amalg *
	\]
	since the terminal object $*$ is a unit for $\times $. Universality of coproduct gives an arrow
	\[
		v_+\amalg w_+ \to v\amalg w\amalg*
	\]
	which the canonical arrow $v_+\amalg w_+\to v_+\times w_+$ factors through by uniqueness. Hence we obtain a diagram
	\[\begin{tikzcd}[row sep = 2.5em, column sep = 2.5em]
	v_+\amalg w_+\ar[two heads, r] \ar[d]  & v\amalg w\amalg *\ar[r]  & v\times w\amalg v\amalg w\amalg *\ar[d]  \\
	*\ar[r]  & *\ar[r]  & v\times w\amalg *=(v\times w)_+
	\end{tikzcd}\]
	since the canonical map doesn't act on $v\times w$, which is a double pushout, since the top left arrow is epi.
\end{proof}

A convenient category $\Top $ of spaces gives a closed symmetric monoidal category $(\Top _*,\wedge,S^0)$. So does $(\sSet_*,\wedge,\partial \Delta^1)$.

\section{Underlying categories of enriched categories}

Here is some motivation which I'm copying down less for the motivation, and more because I think the new concepts are cool. Fix a group $G$ to which the discrete topology has been associated, and take $\Top $ some well-behaved category of topological spaces. The functor category $\Top ^G$ has as objects $G$-spaces and as morphisms $G$-equivariant maps. Unpacking this quickly, a functor is a group action of $G$ on a topological space, meaning a choice of continuous automorphisms of a space that respect the group operation. Morphisms of $G$-spaces are natural transformations, but since grououps are singletons, this amounts to a single continuous map that preserves the group action, called $G$-equivariant.

This category admits two natural enrichments. First, the set of $G$-equivariant maps can be viewed as a subspace $\underline{\Top }^G(X,Y)\subseteq \underline{\Top }(X,Y)$, where the latter has the \emph{compact-open topology} generated by sets
\[
	V(K,U)\coloneqq \left\{ f\in \Top (X,Y) \mid f(K)\subseteq U \right\} ,\qquad K\subseteq X\text{ compact} ,\qquad U\subseteq Y\text{ open} 
.\]
This gives a topological enrichment.

Additionally, $\Top ^G$ is symmetric cartesian monoidal, with product given by the standard cartesian product together with the diagonal action $(g,h)(x,y)=(gx,hy)$. The singleton with the trivial action is a unit thereof. The space $\underline{\Top }(X,Y)$ has a canonical $G$-action
\[
	g\cdot f(x)=gf(g^{-1} x)
.\]
This gives an enrichment over $\Top ^G$, and we write $\underline{\Top }_G(X,Y)$ for this space to distinguish from the other enrichment. With this definition $\Top ^G$ is cartesian closed.

We think of $\Top ^G$ as the underlying (unenriched) category for both categories, but this intuition is difficult to formalize in full generality. However, generality in this case helps extract the full definition.

To define an underlying category, we first need a way to construct underlying sets, and hence a functor $\mathcal{V}\to \Set $. For example, in $\Ab$, $\sSet$, $\Top $, $\Cat $, and more, the underlying set functor is simply the representable functor $\mathcal{V} (*,-): \mathcal{V}  \to \Set $, since $*$ is a singleton. Since this is the unit for the cartesian monoidal structures of these categories, we might generalize this to any symmetric monoidal category with $*$ as the unit. This is on the right track, but we need a bit more.
\begin{definition}[Lax Monoidal]\label{dfn:3.3}
	A functor $F : \mathcal{V}  \to \mathcal{U} $ between symmetric monoidal categories $(\mathcal{V} ,\otimes ,*)$ and $(\mathcal{U} , \Box,1)$ is \emph{lax monoidal} if there are natural transformations (NOT ISOMORPHISMS)
	\[
		F(v)\otimes F(v')\to F(v\,\Box \,v'),\qquad 1\to F(*)
	.\]
\end{definition}
Clearly all strong monoidal functors are lax, with natural transformations as isomorphisms. 

\begin{lemma}[Change of Base]\label{lma:3.2}
	Let $F : \mathcal{V}  \to \mathcal{U} $ be lax monoidal. Then any $\mathcal{V} $-category has an associated $\mathcal{U} $-category with the same objects.
\end{lemma}
\begin{proof}
	Let $\underline{\mathcal{D} }_\mathcal{V} $ be a $\mathcal{V} $-category. Define a $\mathcal{U} $-category $\underline{\mathcal{D} } _\mathcal{U} $ with the same objects, and with hom-objects $\underline{\mathcal{D} } _\mathcal{U}(x,y) \coloneqq F\underline{\mathcal{D} } _\mathcal{V} (x,y)$. The lax monoidal structure of $F$ gives composition and unit maps
	\[
		F\underline{\mathcal{D} } _\mathcal{V} (y,z)\,\Box\, F\underline{\mathcal{D} } _\mathcal{V} (x,y)\to F\left( \underline{\mathcal{D} } _\mathcal{V} (y,z)\otimes \underline{\mathcal{D} }_\mathcal{V} (x,y) \right) \xlongrightarrow{F(\circ )}F\underline{\mathcal{D} } _\mathcal{V} (x,z)=\underline{\mathcal{D} } _\mathcal{U} (x,z);
	\]
	\[
		1\to F(*)\xlongrightarrow{F(\mathrm{id}_x)} F\underline{\mathcal{D} } _\mathcal{V} (x,x)=\underline{\mathcal{D} } _\mathcal{U} (x,x)
	.\]
\end{proof}

The point is that $\mathcal{V} (*,-)$ is lax monoidal. Let $v,w\in \mathcal{V} $. Then the bifunctor $\otimes $ gives for any pair $*\to v$ and $*\to w$ a map $*\otimes * \to v\otimes w$. Precomposing with the $\otimes $-diagonal map, indeed an isomorphism here, gives rise to a function
\[
	\mathcal{V} (*,v)\times \mathcal{V} (*,w)\to \mathcal{V} (*,v\otimes w)
\]
natural in $v$ and $w$, and hence $\mathcal{V} (*,-)$ is lax functorial. We then apply \cref{lma:3.2} to change the base for the enrichment from $\mathcal{V} $ to $\Set $.

\begin{definition}[Underlying Category]\label{dfn:3.4}
	Given a $\mathcal{V} $-category $\underline{\mathcal{C} } $, the underlying category $\mathcal{C} _0$ has the same objects, and as hom-sets
	\[
		\mathcal{C} _0(x,y)\coloneqq \mathcal{V} (*,\underline{\mathcal{C} } (x,y))
	.\]
	Definitions are as in \cref{lma:3.2}: $1_x=\mathrm{id} _x\in \mathcal{V} (*,\underline{\mathcal{C} } (x,x))$ and
	\[\begin{tikzcd}[row sep = 3em, column sep = 3em]
	\mathcal{C} _0(y,z)\times \mathcal{C} (x,y)\ar[" \circ  ", dashed, rr] \ar[equals, d]  &  & \mathcal{C} _0(x,z)\ar[equals, d]  \\
	\mathcal{V} (*,\underline{\mathcal{C} } (y,z))\times \mathcal{V} (*,\underline{\mathcal{C} } (x,y))\ar[" \otimes  "', r]  & \mathcal{V} (*,\underline{\mathcal{C} } (y,z)\otimes \underline{\mathcal{C} } (x,y))\ar[" \mathcal{V} (*\text{,} \circ ) "', r]  & \mathcal{V} (*,\underline{\mathcal{C} } (x,z)).
	\end{tikzcd}\]
\end{definition}

Reconsiling our previous example, first consider the $\Top $-enriched category $\underline{\Top } ^G$. The hom-sets of the underling category are $\Top (*,\underline{\Top } ^G(X,Y))$, equivalently the set of points in $\underline{\Top } ^G(X,Y)$, equivalently all $G$-equivariant maps $X\to Y$. Hence the underlying category is $\Top ^G$. Similarly, the hom-sets of the underlying category of the $\Top ^G$-enriched category $\underline{\Top } _G$ are given by $\Top ^G(*,\underline{\Top } _G(X,Y))$. Since $*$ is the trivial group action, $G$-equivariant maps $*\to \underline{\Top } _G(X,Y)$ are $G$-fixed points. Note that fixed points for the conjugation action are precisely $G$-equivariant maps:
\[
	g\cdot f(x)=gf(g^{-1} x)=gg^{-1} f(x)=f(x)
.\]
Hence, the underlying category is once again $\Top ^G$.

These considerations with group actions generalize. For $\mathcal{V} $ a complete and cocomplete closed symmetric monoidal category, $\mathcal{V} ^G$ inherits a canonical closed symmetrical monoidal structure such that the forgetful functor $U : \mathcal{V} ^G \to \mathcal{V} $ preserves the unit, tensor, and internal homs. The hom-object $\underline{\mathcal{V} } ^G(c,d)$ agrees with $\underline{\mathcal{V} } (c,d)$, and the underlying category is again the original $\mathcal{V} ^G$. The discussion also generalizes in another way.

\begin{lemma}\label{lma:3.3}
	If $(V,\otimes ,*)$ is a closed symmetric monoidal category, the underlying category of the $\mathcal{V} $-category $\underline{\mathcal{V} } $ is $\mathcal{V} $.
\end{lemma}
\begin{proof}
	Internal homs give the adjunction
	\[
		\mathcal{V} (*,\underline{\mathcal{V} } (x,y))\cong \mathcal{V} (*\otimes x,y)\cong \mathcal{V} (x,y)
	,\]
	with the second isomorphism the unitor. Hence, hom-sets of the underlying category agree with $\mathcal{V} $. We now need to show identities and composition agree. The first one is handled by functorality of $\mathcal{V} (*,-)$, so now we need only show composition agrees.

	Using this above isomorphism, we take $f : x \to y$ in $\mathcal{V} $ and the morphism $f : * \to \underline{\mathcal{V} } (x,y)$ to be equivalent, hence we use the same label. Let $\varepsilon _y^x : \underline{\mathcal{V} } (x,y)\otimes x \to y$ be the relevant counit of the relevant adjunction. Following the isomorphism gives a precise definition of $f : x \to y$ in terms of $f : * \to \underline{\mathcal{V} } (x,y)$ as the composition
	\[\begin{tikzcd}[row sep = 2.5em, column sep = 2.5em]
		x\cong *\otimes x \ar[" f\otimes 1 ", r]  & \underline{\mathcal{V} } (x,y)\otimes x\ar[" \varepsilon _y^x ", r]  & y.
	\end{tikzcd}\]
	Now we follow composition as in \cref{dfn:3.4}. The composition in $\underline{\mathcal{V} } $ is defined as
	\[\begin{tikzcd}[row sep = 2.5em, column sep = 2.5em]
	*\cong *\otimes *\ar[" g\otimes f ", r]  & \underline{\mathcal{V} } (y,z)\otimes \underline{\mathcal{V} } (x,y)\ar[" \circ  ", r]  & \underline{\mathcal{V} } (x,z).
	\end{tikzcd}\]
	Recall composition was adjunct to
	\[\begin{tikzcd}[row sep = 2.5em, column sep = 2.5em]
		\underline{\mathcal{V} } (y,z)\otimes \underline{\mathcal{V} } (x,y)\otimes x\ar[" 1\otimes \varepsilon _y^x ", r]&  \underline{\mathcal{V} } (y,z)\otimes y\ar[" \varepsilon _z^y ", r]  & z .
	\end{tikzcd}\]
	Naturality and precomposition with a unitor gives that our definition of composition in its entirety is adjunct to
	\[\begin{tikzcd}[row sep = 2.5em, column sep = 2.5em]
	*\otimes x\cong *\otimes *\otimes x\ar[" g\otimes f\otimes 1 ", r]  & \underline{\mathcal{V} } (y,z)\otimes \underline{\mathcal{V} } (x,y)\otimes x\ar[" 1\otimes \varepsilon _y^x ", r]  & \underline{\mathcal{V} } (y,z)\otimes y\ar[" \varepsilon _z^y ", r]  & z.
	\end{tikzcd}\]
	Functorality and another unitor gives
	\[\begin{tikzcd}[row sep = 2.5em, column sep = 2.5em]
	x\cong *\otimes *\otimes x\ar[" g\otimes 1\otimes 1 ", r]  & \underline{\mathcal{V} } (y,z)\otimes *\otimes x\ar[" 1\otimes f\otimes 1 ", r]  & \underline{\mathcal{V} } (y,z)\otimes \underline{\mathcal{V} } (x,y)\otimes x\ar["1\otimes  \varepsilon _y^x ", r]  & \underline{\mathcal{V} } (y,z)\otimes y\ar[" \varepsilon _z^y ", r]  & z.
	\end{tikzcd}\]
	By functorality of tensor, we can replace the part of the composition
	\[\begin{tikzcd}[row sep = 2.5em, column sep = 2.5em]
	*\otimes x\ar[" f \otimes 1", r]  & \underline{\mathcal{V} } (x,y)\otimes x\ar[" \varepsilon _y^x ", r]  & y
	\end{tikzcd}\]
	with the definition of $f : x \to y$ in $\mathcal{V} $:
	\[\begin{tikzcd}[row sep = 2.5em, column sep = 2.5em]
	x\cong *\otimes x\ar[" g\otimes 1 ", r]  & \underline{\mathcal{V} } (y,z)\otimes x\ar[" 1\otimes f ", r]  & \underline{\mathcal{V} } (y,z)\otimes y\ar[" \varepsilon _z^y ", r]  & z.
	\end{tikzcd}\]
	Functorality of $\otimes $ gives that this composition is equivalent to
	\[\begin{tikzcd}[row sep = 2.5em, column sep = 2.5em]
	x\ar[" f ", r]  & y\cong *\otimes y\ar[" g\otimes 1 ", r]  & \underline{\mathcal{V} } (y,z)\otimes y\ar[" \varepsilon _z^y ", r]  & z.
	\end{tikzcd}\]
	The tail of this composition is precisely the definition of $g$ as an arrow in the underlying category, hence this is the composition $g\circ f$.
\end{proof}

\begin{exercise}
	Let $\underline{\mathcal{D} } $ be a $\mathcal{V} $-category and $g : * \to \underline{\mathcal{D} } (y,z)$ an arrow in $\mathcal{D} _0$. Define for any $x\in \underline{\mathcal{D} } $ postcomposition:
	\[\begin{tikzcd}[row sep = 2.5em, column sep = 2.5em]
	g_*: \underline{\mathcal{D} } (x,y)\cong *\otimes \underline{\mathcal{D} } (x,y)\ar[" g\otimes 1 ", r]  & \underline{\mathcal{D} } (y,z)\otimes \underline{\mathcal{D} } (x,y)\ar[" \circ ", r]  & \underline{\mathcal{D} } (x,z).
	\end{tikzcd}\]
	Precomposition is defined similarly. This definition gives an (unenriched) functor
	\[
		\underline{\mathcal{D} } (x,-): \mathcal{D} _0 \to \mathcal{V} ,\qquad y\mapsto \underline{\mathcal{D} } (x,y),\qquad g\mapsto g_*
	.\]
	Show that the composition with the underlying set functor $\mathcal{V} (*,-)\circ \underline{\mathcal{D} } (x,-)$ is the representable functor $\mathcal{D} _0(x,-): \mathcal{D} _0 \to \Set $.
\end{exercise}
\begin{proof}
	By \cref{dfn:3.4}, $\mathcal{V} (*,\underline{\mathcal{D} } (x,-))\eqqcolon \mathcal{D} _0(x,-)$. Hence, the definitions agree on objects. We also have that $\mathcal{D} _0(x,g)=g_*$ in the conventional sense of the notation, and by definition $\underline{\mathcal{D} } (x,g)=g_*$ in the sense of the notation as above. We also have that for a morphism $f$ in $\mathcal{V} $, $\mathcal{V} (x,f)$ is also precomposition in the conventional sense. Note that $\mathcal{V} (*,\underline{\mathcal{D} } (x,g))=\mathcal{V} (*,g_*)$.

	Let $f : x \to y$ in $\mathcal{D} _0$, and equivalently $f : * \to \underline{\mathcal{D} } (x,y)$ in $\mathcal{V} $. We must show that evaluating $\mathcal{V} (*,g_*)(f)$ is precisely the same as $\mathcal{D} _0(x,g)(f)$. Since $g_*$ is defined as an arrow in $\mathcal{V} $, the former is given as $g_*\circ f$ with composition as defined for $\mathcal{V} $. The latter is given as $g\circ f$, with composition as defined in $\mathcal{D} _0$. Writing out the composition $g_*\circ f$ as defined above gives the composition
	\[\begin{tikzcd}[row sep = 3em, column sep = 3em]
	*\ar[" f ", r]  & \underline{\mathcal{D} } (x,y)\cong *\otimes \underline{\mathcal{D} } (x,y)\ar[" g\otimes 1 ", r]  & \underline{\mathcal{D} } (y,z)\otimes \underline{\mathcal{D} } (x,y)\ar[" \circ  ", r]  & \underline{\mathcal{D} } (x,z)
	\end{tikzcd}\]
	with composition as in $\underline{\mathcal{D} } $. Functorality simplifies this to
	\[\begin{tikzcd}[row sep = 3em, column sep = 2.5em]
	*\cong *\otimes *\ar[" g\otimes f ", r]  & \underline{\mathcal{D} } (y,z)\otimes \underline{\mathcal{D} } (x,y)\ar[" \circ  ", r]  & \underline{\mathcal{D} } (x,z).
	\end{tikzcd}\]
	By \cref{dfn:3.4}, this is precisely the same as the definition of $g\circ f$ in $\mathcal{D} _0$. Thus, $\mathcal{V} (*,\underline{\mathcal{D} } (x,-))=\mathcal{D} _0(x,-)$.
\end{proof}

Hence we have a good notion of precomposition and postcomposition. Importantly, pre- and postcomposition are continuous in topological categories, but we can indeed show more. These representable functors are $\mathcal{V} $-functors.

\section{Enriched functors and enriched natural transformations}

Small $\mathcal{V} $-categories can form a 2-category (if $\mathcal{V} $ is allowed to be symmetric). To make precise this notion we require the notion of a $\mathcal{V} $-functor.

\begin{definition}[$\mathcal{V}$-Functor]\label{dfn:3.5}
	A $\mathcal{V} $-functor $F : \underline{\mathcal{C} }  \to \underline{\mathcal{D} } $ between $\mathcal{V} $-categories is an object map $x\mapsto F(x)$ from $\underline{\mathcal{C} } $ to $\underline{\mathcal{D} } $, together with morphisms
	\[
		F_{x,y} : \underline{\mathcal{C} } (x,y) \to \underline{\mathcal{D} } (x,y)
	\]
	in $\mathcal{V} $ for each pair $x,y\in \underline{\mathcal{C} } $, such that the following diagrams commute for all $x,y,z\in \underline{\mathcal{C} } $:
	\[\begin{tikzcd}[row sep = 3.5em, column sep = 3.5em]
	\underline{\mathcal{C} } (y,z)\otimes \underline{\mathcal{C} } (x,y)\ar[" \circ  ", r] \ar[" F_{y\text{,} z} \otimes F_{x\text{,} y}  "', d]  & \underline{\mathcal{C} } (x,y)\ar[" F_{x\text{,} z}  ", d]  & *\ar[" \mathrm{id} _x ", r] \ar[" \mathrm{id} _{F(x)}  "', dr]  & \underline{\mathcal{C} } (x,x)\ar[" F_{x\text{,} x}  ", d] \\
	\underline{\mathcal{D} } (F(y),F(z))\otimes \underline{\mathcal{D} } (F(x),F(y))\ar[" \circ  "', r]  & \underline{\mathcal{D} } (F(x),F(z)) &  & \underline{\mathcal{D} } (F(x),F(x)).
	\end{tikzcd}\]
\end{definition}

For a regular category, hence a $\Set $-category, these diagrams demand that $F(-\circ -)=F(-)\circ F(-)$ and $F(1_{(-)} )=1_{F(-)} $. More importantly, a topological functor between $\Top $-categories consists of continuous maps between hom-spaces. $\Top $-functors are frequently called \emph{continuous} functors in the literature, an unfortunate conflict of terminology from the uncommon, but still used, notion of a continuous functor being one which preserves small limits.

A simplicial functor between $\sSet$-categories consists of maps between the $n$-simplices of the appropriate hom-spaces. A 2-functor between $\Cat $-categories, equivalently 2-categories, is a map of objects, morphisms, and 2-cells that preserves domains, codomains, composition, and identities at all levels. Going back to a definition we explored earlier, a covariant $\Ab$-functor from a singleton $\Ab$-category $R$ to $\underline{\Ab} $ is a left $R$-module; a contravariant $\Ab$-functor is a right $R$-module.

\begin{exercise}
	Extend \cref{dfn:3.4} to define the underlying functor of $\mathcal{V} $-functors and show that your definition is functorial, i.e., defines a functor $(-)_0: \mathcal{V} \text{-}\Cat  \to \Cat $.
\end{exercise}
\begin{proof}
	Given a $\mathcal{V} $-functor $F : \underline{\mathcal{C} }  \to \underline{\mathcal{D} } $, define $F_0: \mathcal{C} _0 \to \mathcal{D} _0$ to be the same as $F$ on objects, and on morphisms is simply given by the relevant postcomposition in $\mathcal{V} $. More precisely, for $f : x \to y$ in $\mathcal{C} _0$, we have $f : * \to \underline{\mathcal{C} } (x,y)$ in $\mathcal{V} $, and we then define
	\[\begin{tikzcd}[row sep = 2.5em, column sep = 2.5em]
	F_0(f)\coloneqq *\ar[" f ", r]  & \underline{\mathcal{C} } (x,y)\ar[" F_{x\text{,} y}  ", r]  & \underline{\mathcal{D} } (x,y).
	\end{tikzcd}\]
	Hence $F(f)\in \mathcal{V} (*,\underline{\mathcal{D} } (x,y))\eqqcolon \mathcal{D} _0(x,y)$. The fact that $F$ commutes with $\circ $ in $\mathcal{V}$ and that it preserves identities makes it functorial. These are fairly easy to check, but since this is an exercise I will spell it out in completeness. Let $g : y \to z$ and $f : x \to y$ in $\mathcal{C} _0$. We obtain the diagram
	\[\begin{tikzcd}[row sep = 3.5em, column sep = 3.5em]
		*\otimes *\ar[" g\otimes f ", r] &\underline{\mathcal{C} } (y,z)\otimes \underline{\mathcal{C} } (x,y)\ar[" \circ  ", r] \ar[" F_{y\text{,} z} \otimes F_{x\text{,} y}  "', d]  & \underline{\mathcal{C} } (x,y)\ar[" F_{x\text{,} z}  ", d]  \\
						 &\underline{\mathcal{D} } (F(y),F(z))\otimes \underline{\mathcal{D} } (F(x),F(y))\ar[" \circ  "', r]  & \underline{\mathcal{D} } (F(x),F(z)) .
	\end{tikzcd}\]
	This expresses the equality
	\[
		\circ_\mathcal{D}  \circ (F_{y,z} \otimes F_{x,y} )\circ (g\otimes f)=F_{x,z} \circ \circ_\mathcal{C}  \circ (g\otimes f)
	.\]
	Functorality of $\otimes $ gives on the left
	\[
		\circ_\mathcal{D}  \circ  (F_{y,z} \otimes F_{x,y} )\circ  (g\otimes f)=\circ_\mathcal{D} \circ ((F_{y,z} \circ g)\otimes (F_{x,y} \circ f))
	.\]
	We thus obtain by definition of $F_0$ and of composition in $\mathcal{C} _0$ that
	\[
		F_0(g)\circ F_0(f)=F_0(g\circ f)
	.\]
	The identity is manifestly preserved. Thus, $F_0$ is functorial.

	Now for the functor $(-)_0$. We wish to show the assignment
	\[
		\underline{\mathcal{C} } \mapsto \mathcal{C} _0,\qquad F\mapsto F_0
	\]
	is functorial. This is fairly simple. Let $F : \underline{\mathcal{C} }  \to \underline{\mathcal{D} } $ and $G : \underline{\mathcal{D} }  \to \underline{\mathcal{E} } $ be $\mathcal{V} $-functors. Then by definition of functor composition, we have
	\[
		(G\circ F)_0(f)=(G\circ F)_{x,y} \circ f = G_{x,y} \circ F_{x,y} \circ f=(G_0\circ F_0)(f)
	.\]
	Thus the assignment is indeed functorial.
\end{proof}

When $\underline{\mathcal{V} } $ is a closed $\mathcal{V} $-category, there is for any $\underline{\mathcal{C} } $ a \emph{representable functor} $\underline{\mathcal{C} } (c,-): \underline{\mathcal{C} }  \to \underline{\mathcal{V} } $ whose underlying functor is the representable functor $\underline{\mathcal{C} } (c,-): \mathcal{C} _0 \to \mathcal{V} $ as constructed earlier. The object map is merely $x\mapsto \underline{\mathcal{C} } (c,x)$; the morphisms
\[
	\underline{\mathcal{C} } (c,-)_{x,y} : \underline{\mathcal{C} } (x,y) \to \underline{\mathcal{V} } (\underline{\mathcal{C} } (c,x),\underline{\mathcal{C} } (c,y))
\]
are the adjuncts of composition. More precisely, since $\underline{\mathcal{V} } $ is closed, composition
\[
	\circ : \underline{\mathcal{C} } (x,y)\otimes \underline{\mathcal{C} } (c,x) \to \underline{\mathcal{C} }(c,y)
\]
has an adjunct under \cref{dfn:3.2} to
\[
	\underline{\mathcal{C} } (x,y)\to \underline{\mathcal{V} } (\underline{\mathcal{C} } (c,x),\underline{\mathcal{C} } (c,y))
.\]
The definition of composition in the underlying category recovers $g\mapsto g_*$ for the underlying functor.

Suppose still that $\mathcal{V} $ is closed monoidal. Recall an unenriched category $\mathcal{C} $ admits a free $\mathcal{V} $-categorical structure with objects the same as $\mathcal{C} $, and hom-objects given by
\[
	\underline{\mathcal{C} } (a,b)\coloneqq \mathcal{C} (a,b)\cdot * = \coprod_{\mathcal{C} (a,b)} *
.\]
A functor $F$ from the free $\mathcal{V} $-category on $\mathcal{C} $ to a $\mathcal{V} $-category $\underline{\mathcal{D} } $ consists of an object function $c\mapsto F(c)$, together with morphisms
\[
	F_{x,y} :\coprod_{\mathcal{C} (x,y)} * \to \underline{\mathcal{D} } (F(x),F(y))
.\]
Universality of coproduct gives that $F_{x,y} $ is completely determined by some family of morphisms $X\coloneqq \left\{ \varphi_f : * \to \underline{\mathcal{D} } (F(x),F(y)) \right\}_{f\in \mathcal{C} (x,y)} $. Since
\[
	X\subseteq \mathcal{V} (*,\underline{\mathcal{D} } (F(x),F(y)))\eqqcolon \mathcal{D} _0(F(x),F(y)),
,\]
we have that $F_{x,y} $ is equivalently an assignment $\mathcal{C} (x,y)\to \mathcal{D} _0(F(x),F(y))$. Assigning the notation $F^V(\mathcal{C} )$ to the free $\mathcal{V} $-category on $\mathcal{C} $ (this notation will never be used again and is simply here to write down precisely my point), we have that functors $F : F^V(\mathcal{C} ) \to \underline{\mathcal{D} } $ are precisely functors $F : \mathcal{C}  \to \mathcal{D} _0$. Hence, there is an isomorphism
\[
	\mathcal{V} \text{-} \Cat (F^V(\mathcal{C} ),\underline{\mathcal{D} } )\cong \Cat (\mathcal{C} ,\mathcal{D} _0)
\]
natural in both entries by universality of coproduct. Thus there is an adjunction $F^V\dashv (-)_0$, and for this reason we will never again use the notation $F^V$, since functors $F^V(\mathcal{C} )\to \underline{\mathcal{D} } $ are simply functors $\mathcal{C} \to \mathcal{D} _0$.

Recall that $\sSet$ admits a closed enrichment via $\underline{\sSet} (X,Y)_n=\sSet(\Delta^n\times X,Y)$. The terminal simplicial set $\Delta^0 = \boldsymbol{\Delta} (-,[0])$ in an unenriched sense also represents the functor $(-): \sSet \to \Set $, which takes a simplicial set $X$ to it's set $X_0$ of vertices. However viewing $\sSet$ as self-enriched, the represented functor $\underline{\sSet} (\Delta^0,-)$ becomes the identity on $\sSet$:
\[
	\underline{\sSet} (\Delta^0,X)_n=\sSet(\Delta^n \times \Delta^0,X)\cong\sSet(\Delta^n,X)\cong X_n
.\]
The equality is definition, the first isomorphism is that $\Delta^0$ is terminal, and the second is Yoneda.

The next natural definition is that of a $\mathcal{V} $-natural transformation. However, we need to recall that there are no morphisms to act on, so the natural way to reproduce the naturality square
\[\begin{tikzcd}[row sep = 2.5em, column sep = 2.5em]
F(x)\ar[" F(f) ", r] \ar[" \alpha_x "', d]  & F(y)\ar[" \alpha_y ", d]  \\
G(x)\ar[" G(f) "', r]  & G(y)
\end{tikzcd}\]
is via pre- and post-composition with components of $\alpha$.

\begin{definition}[$\mathcal{V} $-Natural Transformation]\label{dfn:3.6}
	A $\mathcal{V} $-natural transformation $\alpha : F \Rightarrow G$ between $\mathcal{V} $-functors $F,G : \underline{\mathcal{C} }  \to \underline{\mathcal{D} } $ is a morphism $\alpha_x : * \to \underline{\mathcal{D} } (F(x),G(x))$ in $\mathcal{V} $ for each $x\in \underline{\mathcal{C} } $, such that for all $x,y\in \underline{\mathcal{C} } $, the diagram
	\[\begin{tikzcd}[row sep = 2.5em, column sep = 2.5em]
	\underline{\mathcal{C} } (x,y)\ar[" F_{x\text{,} y}  ", r] \ar[" G_{x\text{,} y}  "', d]  & \underline{\mathcal{D} } (F(x),F(y))\ar[" (\alpha_y)_* ", d]  \\
	\underline{\mathcal{D} } (G(x),G(y))\ar[" (\alpha_x)^* "', r]  & \underline{\mathcal{D} } (F(x),G(y))
	\end{tikzcd}\]
	with precomposition $(\alpha_x)^*$ and postcomposition $(\alpha_y)_*$ defined as earlier.
\end{definition}

The intuitive understanding of this in unenriched categories is that on the top right path, morphisms $f : x \to y$ in $\mathcal{C} $ should be carried to morphisms $F(f): F(x) \to F(y)$ in $\mathcal{D} $, and then postcomposed with $\alpha_y$. The bottom left path says that this shouldbe equal to taking your morphism and applying $G$ to obtain $G(f): G(x) \to G(y)$, then precompositing with $\alpha_x$, giving an equality
\[
	\alpha_y \circ F(f)=G(f)\circ \alpha_x
.\]
This is precisely the naturality condition.

\begin{example}
	An $\Ab$-natural transformation between two left $R$-modules $A,B : R \to \Ab$ is a single arrow $\alpha_* : * \to \Ab(A,B)$ in $\Ab$, and hence a group homomorphism, that must commute with scalar multiplication by our above discussion. Hence $\Ab$-natural transformations here are $R$-module homomorphisms, and thus $\Ab\text{-} \Fun(R,\Ab)\cong\Rmod_R$.
\end{example}

The addition of $\mathcal{V} $-natural transformations assembles $\mathcal{V} \text{-} \Cat $ into a 2-category. We've shown already that the underlying set functor $(-)_0$ extends to $\mathcal{V} $-functors, making it a functor $\mathcal{V} \text{-} \Cat \to \Cat $. We now claim we can use it to define underlying $\mathcal{V} $-natural transformations.

\begin{proposition}\label{prp:3.1}
	Underlying $(-)_0$ is a $2$-functor $\mathcal{V} \text{-} \Cat \to \Cat $. Additionally, lax monoidal functors $F : \mathcal{V}  \to \mathcal{U}$ induce 2-functors $\mathcal{V} \text{-} \Cat \to \mathcal{U} \text{-} \Cat$.
\end{proposition}
\begin{proof}
	We prove the second statement first. By \cref{lma:3.2}, we already have an object function by mapping some $\mathcal{V} $-category $\underline{\mathcal{D} } _\mathcal{V} $ to the $\mathcal{U} $-category $\underline{\mathcal{D} } _\mathcal{U} $ with the same objects, and hom-objects
	\[
		\underline{\mathcal{D} } _\mathcal{U} (x,y)\coloneqq F\left( \underline{\mathcal{D} } _\mathcal{V} (x,y) \right) 
	.\]
	Let $G : \underline{\mathcal{C} }_\mathcal{V}  \to \underline{\mathcal{D} } _\mathcal{V} $ be a $\mathcal{V} $-functor. Define a new functor $F(G)$ by $F(G)=G$ on objects, and
	\[
		(F(G))_{x,y} =F(G_{x,y} )
	.\]
	Since $G$ is $\mathcal{V} $-functorial, the diagram
	\[\begin{tikzcd}[row sep = 3.5em, column sep = 3.5em]
	\underline{\mathcal{C} } _\mathcal{V} (y,z)\otimes \underline{\mathcal{C} } _\mathcal{V} (x,y)\ar[" \circ  ", r] \ar[" G_{y\text{,} z} \otimes G_{x\text{,} y}  "', d]  & \underline{\mathcal{C} } _\mathcal{V} (x,z)\ar[" G_{x\text{,} z}  ", d]  \\
	\underline{\mathcal{D} } _\mathcal{V} (y,z)\otimes \underline{\mathcal{D} } _\mathcal{V} (x,y)\ar[" \circ  "', r]  & \underline{\mathcal{D} } _\mathcal{V} (x,z)
	\end{tikzcd}\]
	commutes in $\mathcal{V} $. Hence by functorality of $F$, the diagram
	\[\begin{tikzcd}[row sep = 3.5em, column sep = 3.5em]
	F(\underline{\mathcal{C} } _\mathcal{V} (y,z)\otimes \underline{\mathcal{C} } _\mathcal{V} (x,y))\ar[" F(\circ ) ", r] \ar[" F(G_{y\text{,} z} \otimes G_{x\text{,} y} ) "', d]  & \underline{\mathcal{C} } _\mathcal{U} (x,z)\ar[" F(G_{x\text{,} z} ) ", d]  \\
	F(\underline{\mathcal{D} } _\mathcal{V} (y,z)\otimes \underline{\mathcal{D} } _\mathcal{V} (x,y))\ar[" F(\circ ) "', r]  & \underline{\mathcal{D} } _\mathcal{U} (x,z)
	\end{tikzcd}\]
	commutes in $\mathcal{U} $, where we have swapped for the $\mathcal{U} $-categories where possible. Recall since $F$ is lax monoidal, there is a natural transformation
	\[
		F(v)\otimes F(v')\to F(v\otimes v')
	.\]
	More precisely, there is a natural transformation
	\[
		F(-)\otimes F(-)\Rightarrow F(-\otimes -)
	.\]
	Denote this natural transformation by $\lambda$. Since precomposing in the corner preserves commutivity, the diagram
	\[\begin{tikzcd}[row sep = 3.5em, column sep = 3.5em]
	\underline{\mathcal{C} } _\mathcal{U} (y,z)\otimes \underline{\mathcal{C} } _\mathcal{U} (x,y)\ar[" \lambda ", r] & F(\underline{\mathcal{C} } _\mathcal{V} (y,z)\otimes \underline{\mathcal{C} } _\mathcal{V} (x,y))\ar[" F(\circ ) ", r] \ar[" F(G_{y\text{,} z} \otimes G_{x\text{,} y} ) "', d]  & \underline{\mathcal{C} } _\mathcal{U} (x,z)\ar[" F(G_{x\text{,} z} ) ", d]  \\
	& F(\underline{\mathcal{D} } _\mathcal{V} (y,z)\otimes \underline{\mathcal{D} } _\mathcal{V} (x,y))\ar[" F(\circ ) "', r]  & \underline{\mathcal{D} } _\mathcal{U} (x,z)
	\end{tikzcd}\]
	commutes in $\mathcal{U} $. Thus, by naturality of $\lambda$, the diagram
	\[\begin{tikzcd}[row sep = 3.5em, column sep = 3.5em]
	\underline{\mathcal{C} } _\mathcal{U} (y,z)\otimes \underline{\mathcal{C} } _\mathcal{U} (x,y)\ar[" \lambda ", r] \ar[" F(G_{y\text{,} z} )\otimes F(G_{x\text{,} y} ) "', d] & F(\underline{\mathcal{C} } _\mathcal{V} (y,z)\otimes \underline{\mathcal{C} } _\mathcal{V} (x,y))\ar[" F(\circ ) ", r] \ar[" F(G_{y\text{,} z} \otimes G_{x\text{,} y} ) "', d]  & \underline{\mathcal{C} } _\mathcal{U} (x,z)\ar[" F(G_{x\text{,} z} ) ", d]  \\
	\underline{\mathcal{D} } _\mathcal{U} (y,z)\otimes \underline{\mathcal{D} } _\mathcal{U} (x,y)\ar[" \lambda "', r] & F(\underline{\mathcal{D} } _\mathcal{V} (y,z)\otimes \underline{\mathcal{D} } _\mathcal{V} (x,y))\ar[" F(\circ ) "', r]  & \underline{\mathcal{D} } _\mathcal{U} (x,z)
	\end{tikzcd}\]
	commutes as well. Recall by \cref{lma:3.2} that composition is given in the $\mathcal{U} $-categories via
	\[\begin{tikzcd}[row sep = 2.5em, column sep = 3.5em]
		\underline{\mathcal{C} } _\mathcal{U} (y,z)\otimes \underline{\mathcal{C} } _\mathcal{U} (x,y)\ar[" \lambda ", r] & F\left( \underline{\mathcal{C} } _\mathcal{V} (y,z)\otimes \underline{\mathcal{C} } _\mathcal{V} (x,y) \right) \ar[" F(\circ  )", r] & \underline{\mathcal{C} } _\mathcal{U} (x,z).
	\end{tikzcd}\]
	Note that this is the top and bottom compositions of the above diagram, so simplifying we obtain the diagram
	\[\begin{tikzcd}[row sep = 3.5em, column sep = 3.5em]
	\underline{\mathcal{C} } _\mathcal{U} (y,z)\otimes \underline{\mathcal{C} } _\mathcal{U} (x,y)\ar[" F(G_{y\text{,} z} )\otimes F(G_{x\text{,} y} ) "', d] \ar[" \circ  ", r]  & \underline{\mathcal{C} }_\mathcal{U} (x,z)\ar[" F(G_{x\text{,} z} ) ", d]   \\
	\underline{\mathcal{D} } _\mathcal{U} (y,z)\otimes \underline{\mathcal{D} } _\mathcal{U} (x,y)\ar[" \circ  "', r]  & \underline{\mathcal{D} } _\mathcal{U} (x,z).
	\end{tikzcd}\]
	Thus, $F(G)$ preserves composition.

	To show $F(G)$ preserves identities, recall by \cref{dfn:3.5} that since $G$ is $\mathcal{V} $-functorial, for any $x$ the diagram
	\[\begin{tikzcd}[row sep = 3.5em, column sep = 3.5em]
	*\ar[" \mathrm{id} _x ", r] \ar[" \mathrm{id} _{G\left(x \right) }  "', dr]  & \underline{\mathcal{C} } _\mathcal{V} (x,x)\ar[" G_{x\text{,} x}  ", d]  \\
	 & \underline{\mathcal{D} } _\mathcal{V} (G(x),G(x))
	\end{tikzcd}\]
	commutes. Functorality of $F$ gives that the diagram
	\[\begin{tikzcd}[row sep = 3.5em, column sep = 3.5em]
	F(*)\ar[" F(\mathrm{id} _x) ", r] \ar[" F(\mathrm{id} _{G(x)}  )"', dr]  & \underline{\mathcal{C} } _\mathcal{U} (x,x)\ar[" F(G_{x\text{,} x} ) ", d]  \\
	 & \underline{\mathcal{D} } _\mathcal{U} (G(x),G(x))
	\end{tikzcd}\]
	commutes. To avoid confusion, denote by $1$ the unit in $\mathcal{U} $. Since $F$ is lax monoidal, there is a natural transformation $1\to F(*)$, giving that
	\[\begin{tikzcd}[row sep = 3.5em, column sep = 3.5em]
		1\ar[r] & F(*)\ar[" F(\mathrm{id} _x) ", r] \ar[" F(\mathrm{id} _{G(x)}  )"', dr]  & \underline{\mathcal{C} } _\mathcal{U} (x,x)\ar[" F(G_{x\text{,} x} ) ", d]  \\
			& & \underline{\mathcal{D} } _\mathcal{U} (G(x),G(x))
	\end{tikzcd}\]
	commutes. Recall also that \cref{lma:3.2} defines the identity as the composition
	\[\begin{tikzcd}[row sep = 2.5em, column sep = 3.5em]
	\mathrm{id} _x : 1\ar[r]  & F(*)\ar[" F(\mathrm{id} _x) ", r]  & \underline{\mathcal{D} } _\mathcal{U} (x,x).
	\end{tikzcd}\]
	Note this is the top composition on the above diagram, and thus it suffices to show that $F(\mathrm{id} _{G(x)} )\circ (1\to F(*))=\mathrm{id} _{F(G(x))} $. Note that $F(G(x))=G(x)$, and thus $\mathrm{id} _{F(G(x))} =\mathrm{id} _{G(x)} $. It thus suffices to show $F$ fixes identities. Recall from \cref{dfn:3.1} that the identity must satisfy coherence conditions
	\[\begin{tikzcd}[row sep = 3.5em, column sep = 3.5em]
		\underline{\mathcal{C} }_\mathcal{V} (x,y)\otimes *\ar[" 1\otimes \mathrm{id}_x ", r] \ar[" \cong "', dr]  & \underline{\mathcal{C} }_\mathcal{V} (x,y)\otimes \underline{\mathcal{C} }_\mathcal{V} (x,x)\ar[" \circ  ", d]  & \underline{\mathcal{C} }_\mathcal{V} (y,y)\otimes \underline{\mathcal{C} }_\mathcal{V} (x,y)\ar[" \circ  ", d]  & *\otimes \underline{\mathcal{C} }_\mathcal{V} (x,y)\ar[" \mathrm{id} _y \otimes 1 "', l] \ar[" \cong ", dl]  \\
													      & \underline{\mathcal{C} }_\mathcal{V} (x,y) & \underline{\mathcal{c} }_\mathcal{V} (x,y) & 
	\end{tikzcd}\]
	and the same for $\underline{\mathcal{D} } _\mathcal{V} $. The proof will be the same for both diagrams, so we prove only the left diagram. Functorality of $F$ gives that the diagram
	\[\begin{tikzcd}[row sep = 3.5em, column sep = 3.5em]
	F\left( \underline{\mathcal{C} } _\mathcal{V} (x,y)\otimes * \right) \ar[" \cong "', dr] \ar[" F(1\otimes \mathrm{id} _x) ", r]  & F\left( \underline{\mathcal{C} } _\mathcal{V} (x,y)\otimes \underline{\mathcal{C} } _\mathcal{V} (x,x) \right) \ar[" F(\circ)  ", d]  \\
	 & \underline{\mathcal{C} } _\mathcal{U} (x,y)
	\end{tikzcd}\]
	commutes. By naturality of $\lambda$ and the fact that $F(1)=1$, it follows that
	\[\begin{tikzcd}[row sep = 3.5em, column sep = 3.5em]
	\underline{\mathcal{C} } _\mathcal{U}(x,y)\otimes F(*)\ar[" 1\otimes F(\mathrm{id} _x) ", r] \ar[" \lambda "', d]  & \underline{\mathcal{C} } _\mathcal{U} (x,y)\otimes \underline{\mathcal{C} } _\mathcal{U} (x,x)\ar[" \lambda ", d]  \\
	F\left( \underline{\mathcal{C} } _\mathcal{V} (x,y)\otimes * \right) \ar[" F(1\otimes \mathrm{id} _x) ", r] \ar[" \cong "', dr]  & F\left( \underline{\mathcal{C} } _\mathcal{V} (x,y)\otimes \underline{\mathcal{C} } _\mathcal{V} (x,x) \right) \ar[" F(\circ ) ", d]  \\
	 & \underline{\mathcal{C} } _\mathcal{U} (x,y)
	\end{tikzcd}\]
	commutes. am on the wrong track, apply coherence diagrams before naturality to try to get rid of the stupid
\end{proof}

